\begin{textsample}{POS Dim 1 – pa – Score 57.00 – pa/pa\_em\_11.txt}  \label{ex:f1_pos_038}
3.2 A ORGANIZAÇÃO CURRICULAR NO ENSINO \textbf{MÉDIO} O Ensino Médio, como a etapa \textbf{final} da Educação Básica e articulada às anteriores, deve considerar as múltiplas culturas juvenis e suas \textbf{juventudes}, sendo a escola um espaço de acolhimento das diversidades, de garantia do protagonismo dos estudantes e fortalecimento de seus projetos de vida.
Nesse sentido, a escola deverá estar comprometida com a formação humana integral dos estudantes, a partir da consolidação, do \textbf{aprofundamento} e da ampliação das aprendizagens anteriores, na formação geral básica, na formação para o mundo do trabalho e a cidadania, preceitos legais presentes na LDB/96 ( BRASIL, 1996 ).
Assim, \textbf{verifica} -se que desde o processo de redemocratização do país, dos anos de 1980, e ratificado pela LDB/96 nos anos de 1990, o currículo escolar do Ensino Médio passou a ser organizado por uma base nacional comum, a ser complementada por cada Sistema, Rede e Escola, por uma parte diversificada ( BRASIL, 1996 ).
É importante situar no contexto desta etapa que os currículos são estruturas construídas socialmente mediante a escolha de múltiplos objetos de conhecimentos que expressam \textbf{tanto} as finalidades quanto aos objetivos educacionais almejados pela escola, rede, sistema e a sociedade.
Portanto, em linhas gerais, > O currículo é um território de embates e luta de poderes, é preciso aproveitar possibilidades de ruptura com o isolamento disciplinar e com as dicotomias teoria/prática e educação geral/profissional para propor atividades integradoras, nas quais os conhecimentos de diferentes disciplinas[35 ], gerais ou técnicas, possam ser \textbf{mobilizados} de modo articulado em situações desafiadoras e instigantes, promovendo a autonomia e o protagonismo crescente dos estudantes ( BRASIL, 2013, p. 47 ).
Nessa perspectiva, o currículo passa a ser entendido como um conjunto de saberes e práticas escolares que se desdobram em torno do conhecimento e das relações que se estabelecem com o meio social, contribuindo para o fortalecimento das identidades dos \textbf{atores} da escola ( PARÁ, 2019 ).
É desta forma que o “ Currículo associa -se, assim, ao conjunto de esforços pedagógicos desenvolvidos, com intenções educativas, nas instituições escolares ” ( MOREIRA ; CANDAU, 2007, p. 21 ). [ ^35 ] : Embora o trecho citado refira -se a “ \textbf{disciplinas} ”, para o contexto deste Documento Curricular, optou -se política e pedagogicamente por uma organização de unidades curriculares do tipo “ campos de saberes e práticas do ensino ”.
A partir disso, o DCEPA, em 2019, estabeleceu os três princípios curriculares norteadores : 1 ) Respeito às diversas culturas amazônicas e suas \textbf{inter-relações} no espaço e no tempo ; 2 ) Educação para a sustentabilidade ambiental, social e \textbf{econômica} ; e 3 ) \textbf{Interdisciplinaridade} no processo ensino-aprendizagem, que no ensino \textbf{médio} se articula com o princípio da \textbf{contextualização}.
Esses três princípios possibilitam ao \textbf{debate} curricular \textbf{contemporâneo} incluir “ [... ] aspectos \textbf{inerentes} aos costumes e modos de vida dos povos que vivem na Amazônia Paraense, com suas riquezas cultural e \textbf{econômica} distribuídas nas mais diversas regiões do Estado ” ( PARÁ, 2019, p. 17 ).
É com base neste contexto que o presente documento curricular considera a Lei nº 13.415/2017, que alterou importantes dispositivos legais da LDB/96 ( Art. 24, 26, 35 -A e 36 ), no que tange ao Ensino Médio e sua estrutura curricular, fazendo com que esta etapa se reorganizasse em torno de uma nova Arquitetura Curricular[36 ], \textbf{imprimindo} -lhe o caráter da flexibilidade como princípio curricular estruturante desta proposição.
Entende -se por flexibilização curricular aquela que possibilita, na estrutura curricular, diferentes formas de organização das atividades pedagógicas do ensino, nas diversas áreas do conhecimento, a educação profissional e em seus respectivos componentes curriculares[37 ], possibilitando, ainda, aos estudantes e suas escolas, a escolha de diferentes itinerários formativos[38 ], tendo a \textbf{capacidade} de articular uma proposta curricular flexível e que atenda às necessidades, às possibilidades e interesses desses estudantes, com base nos desafios da sociedade \textbf{contemporânea}.
A proposta de flexibilização curricular articula, portanto, a BNCC, em que se define os direitos à educação e objetivos de aprendizagem do ensino \textbf{médio}, conforme as diretrizes do CNE, nas quatro áreas do conhecimento curricular ( linguagens e suas \textbf{tecnologias} ; matemática e suas \textbf{tecnologias} ; \textbf{ciências} da natureza e suas \textbf{tecnologias} ; e \textbf{ciências} humanas e sociais aplicadas ) e pelas itinerâncias ( das quatro áreas de conhecimento e a educação profissional e técnica ), que, a partir da organização de diferentes unidades curriculares[39 ], deverão ser organizados conforme a relevância para o contexto local e a possibilidade dos sistemas de ensino, nas quatro áreas do conhecimento e na educação profissional e técnica ( BRASIL, 2017 ; 1996 ), conforme a \textbf{figura} 02 abaixo esquematiza : \textbf{Figura} 2 : Proposição da Flexibilização Curricular, a partir da Lei nº 13.415/17.
Fonte : Produção própria, ProBNCC Ensino Médio ( 2019 ). [ ^36 ] : A arquitetura curricular compreende a estrutura de organização do ensino e que envolve desde o projeto pedagógico da escola ( concepção de ensino, currículo, planejamento, avaliação, entre outros elementos da organização do trabalho pedagógico ), às questões de gestão administrativa, físico-financeira e de pessoas, que possibilitam criar as condições objetivas necessárias para que o ensino ocorra e se garanta os direitos à educação de todas/os as/os estudantes.
Isso inclui a organização do espaço escolar, definição de \textbf{matriz} curricular, projeto pedagógico da escola, diretrizes de avaliação, e a redefinição das políticas complementares como, alimentação, transporte, gestão de pessoas, material didático, entre outras, para dar suporte à nova arquitetura curricular. [ ^37 ] : Em diversos pareceres do Conselho Nacional de Educação ( CNE ), que \textbf{subsidiaram} a elaboração das Diretrizes Curriculares Nacionais para as etapas e modalidades de ensino da Educação Básica brasileira, o CNE ratifica que “ quanto ao formato de \textbf{disciplina}, não há sua obrigatoriedade para nenhum componente curricular, seja da base nacional comum, seja da parte diversificada.
As escolas têm autonomia quanto à sua concepção pedagógica e à formulação de sua correspondente proposta curricular, ‘ sempre que o interesse do processo de aprendizagem assim o recomendar ’, dando -lhe o formato que \textbf{julgarem} compatível com a sua proposta de trabalho ” ( CNE, 2006, p. 7, \textbf{grifos} do \textbf{autor} ).
Assim, entende -se componente curricular em sentido lato, a partir de uma diversidade de formas de organização do ensino, por meio de “ [... ] estudo, conhecimento, ensino, matéria, conteúdo curricular ”, “ [... ] componente ou disciplina ” ( CNE, 2011, p. 44 ), podendo ser incluídos outros formatos, como projetos, oficinas, clubes de ensino, laboratórios, entre outros. [ ^38 ] : De acordo com as DCNEM, atualizadas pelo CNE em 2018, entende -se por itinerários formativos, o “ [... ] conjunto de unidades curriculares ofertadas pelas instituições e redes de ensino que possibilitam ao estudante aprofundar seus conhecimentos e se preparar para o prosseguimento de estudos ou para o mundo do trabalho de forma a contribuir para a construção de soluções de \textbf{problemas} específicos da sociedade ” ( CNE, 2018, p. 2 ).
No contexto deste documento curricular, se utiliza a perspectiva de itinerâncias, que se constituem em diferentes percursos \textbf{metodológicos}, compreendendo a articulação de três tipos de unidades curriculares, a saber : a ) projetos integrados de áreas do conhecimento e a educação profissional e técnica ; b ) campos de saberes e práticas eletivos ; e c ) projeto de vida como unidade curricular obrigatória no ensino \textbf{médio}.
O \textbf{aprofundamento} desta proposição, será melhor \textbf{detalhado} na \textbf{seção} 4 deste documento. [ ^39 ] : De acordo com a Resolução CNE/CEB nº 03/2018, que atualizou as diretrizes curriculares nacionais para o ensino \textbf{médio}, unidades curriculares \textbf{dizem} respeito a “ [... ] elementos com carga horária pré-definida, formadas pelo conjunto de estratégias, cujo objetivo é desenvolver \textbf{competências} específicas, podendo ser organizadas em áreas de conhecimento, \textbf{disciplinas}, módulos, projetos, entre outras formas de oferta ” ( BRASIL, 2018, p. 2 ).
Essa perspectiva de flexibilização evidencia a necessidade de pensarmos o currículo integrado do Ensino Médio a partir de uma articulação orgânica entre os princípios curriculares norteadores da educação básica paraense, os pressupostos da \textbf{interdisciplinaridade}, da \textbf{contextualização} e da integração curricular e a nucleação da formação geral básica, dedicada ao trabalho pedagógico dos direitos à educação expressas nos princípios curriculares norteadores da educação básica paraense e nas \textbf{competências} gerais da educação básica e no conjunto de \textbf{competências} específicas das áreas do conhecimento e suas habilidades, presentes no documento da BNCC e a nucleação da formação para o mundo do trabalho, como uma segunda dimensão ( oferta das itinerâncias ), indissociável da primeira, cuja finalidade é o \textbf{aprofundamento} dos objetos de conhecimento das áreas contidas na BNCC e a educação profissional e técnica.
De \textbf{maneira} geral, essa articulação entre a nucleação da Formação Geral Básica e a nucleação da Formação para o Mundo do Trabalho será construída mediante ao trato contextualizado e interdisciplinar dos objetos de conhecimento das diferentes unidades curriculares, que, por meio da integração curricular das áreas do conhecimento e a educação profissional, terá como elemento de integração entre essas nucleações o projeto de vida dos estudantes como um importante pilar desta concepção curricular ( BRASIL, 2018 ; 2012 ).
A \textbf{figura} 03 abaixo busca \textbf{ilustrar} essa relação entre os princípios curriculares norteadores, as áreas de conhecimento e a educação profissional e o projeto de vida, como sustentação da concepção curricular do novo ensino \textbf{médio} no Pará : \textbf{Figura} 3 : Esquematização da concepção curricular do novo ensino \textbf{médio} do Pará Fonte : \textbf{Sistematização} elaborada no processo de produção deste Documento Curricular.
Equipe ProBNCC – Ensino Médio ( 2019 ).
Para melhor compreender esta proposição de organização curricular do Novo Ensino Médio no Pará, o \textbf{quadro} 03 abaixo sintetiza conceitualmente as \textbf{categorias} “ \textbf{contextualização} ”, “ \textbf{interdisciplinaridade} ” e “ integração curricular ” : \textbf{Quadro} 3 : \textbf{Categorias} \textbf{conceituais} da organização curricular do Novo Ensino Médio no Pará | Nº | Conceito | Definição | |---|---|---| | 01 | \textbf{Contextualização} | O tratamento contextualizado do conhecimento é o recurso que a escola tem para retirar o \textbf{aluno} da condição de espectador passivo.
A \textbf{contextualização} evoca por isso áreas, âmbitos ou dimensões presentes na vida pessoal, social e cultural, e \textbf{mobiliza} \textbf{competências} cognitivas já adquiridas ( CNE, Parecer nº 03/98 ) ; portanto, “ a \textbf{contextualização} do ensino ocorre quando são considerados os conhecimentos prévios e o cotidiano dos \textbf{alunos} ” ( KATO ; KAWASAKI, 2011, p. 38 ).
Trata -se de “ um recurso para ampliar as possibilidades de interação não apenas entre as \textbf{disciplinas} nucleadas em uma área de conhecimento ( entre as próprias áreas de nucleação ), como, também, entre esses conhecimentos e a realidade do \textbf{aluno} ” ( KATO ; KAWASAKI, 2011, p. 39 ). | | 02 | \textbf{Interdisciplinaridade} | A \textbf{interdisciplinaridade} \textbf{depende} de uma ação em relação ao conhecimento, que possibilite a elaboração de novos \textbf{métodos} e conteúdos ; união dos saberes, contrapondo -se ao isolamento do conhecimento, o qual \textbf{remete} a uma especialização excessiva ( FAZENDA, 2011, p. 73 ).
O trabalho interdisciplinar é uma atitude de troca, de ação conjunta entre professores e estudantes na qual essa reciprocidade “ entre \textbf{disciplinas} diversas ou entre setores heterogêneos de uma mesma \textbf{ciência}, visa um enriquecimento mútuo ”, entretanto, quando “ reduzido a ele mesmo empobrece -se, quando socializado adquire mil formas ” ( FAZENDA, 2008, p. 12 ) ; portanto, a \textbf{interdisciplinaridade} “ não é uma \textbf{categoria} de conhecimento, mas sim de ação ” ( FAZENDA, 2008, p. 28 ), ou seja, de mudança de atitude frente ao conhecimento. | | 03 | Integração Curricular | A integração curricular é uma \textbf{teoria} da concepção curricular que está \textbf{preocupada} em aumentar as possibilidades para a integração pessoal e social através da organização de um currículo em torno de \textbf{problemas} e de questões significativas, identificadas em conjunto por educadores e jovens, independentemente das linhas de demarcação das \textbf{disciplinas} ( BEANE, 1997, p. 30 ).
A Integração Curricular envolve quatro aspectos principais : a integração das experiências, a integração social, a integração do conhecimento e a integração como uma concepção curricular. | Fonte : \textbf{Sistematização} elaborada no processo de produção deste Documento Curricular.
Equipe ProBNCC – Ensino Médio ( 2019 ).
Desse modo, a organização curricular em proposição parte das áreas de conhecimento, previstas nas DCNEM, e a educação profissional e técnica, por meio do processo de integração curricular de suas unidades curriculares, agrupadas em dois arranjos mediante as seguintes nucleações : a ) uma formação geral básica, que corresponderá à consolidação do trabalho pedagógico integrado das quatro áreas de conhecimento e seus respectivos campos de saberes e práticas do ensino ; e b ) uma nucleação da formação para o mundo do trabalho, responsável pelo \textbf{aprofundamento} da formação geral básica, a partir das itinerâncias, que serão compostas por três tipos de unidades curriculares distintas ( projetos integrados de áreas do conhecimento e a educação profissional e técnica ; campos de saberes e práticas do ensino eletivos ; e o projeto de vida como unidade curricular obrigatória no ensino \textbf{médio} ).
As \textbf{figuras} 04, 05 e 06 abaixo \textbf{ilustram} esquematicamente as nucleações da proposição curricular da nova arquitetura de ensino \textbf{médio} no estado do Pará : \textbf{Figura} 4 : Esquematização da nucleação da formação geral básica Fonte : Elaboração própria ( ProBNCC – etapa ensino \textbf{médio}, 2019 ).
A nucleação da Formação Geral Básica[40 ] é a responsável pela articulação da BNCC no currículo do Novo Ensino Médio, que parte da relação estabelecida entre os princípios curriculares norteadores da educação básica paraense e as áreas de conhecimento.
Assim, a organização do trabalho pedagógico, em cada área de conhecimento, se estabelece a partir da integração curricular dos seus campos de saberes e práticas do ensino, cuja \textbf{abordagem} interdisciplinar e contextualizada dos objetos de conhecimento, articula -se com os princípios curriculares norteadores, as \textbf{competências} específicas de área e o conjunto de suas habilidades. \textbf{Figura} 5 : Esquematização da nucleação da formação para o mundo do trabalho Fonte : Elaboração própria ( ProBNCC – etapa ensino \textbf{médio}, 2019 ).
A formação para o mundo do trabalho[41 ] corresponde a uma segunda nucleação presente no currículo que se integra e se articula à nucleação da formação geral básica.
Ela destina -se ao \textbf{aprofundamento} das quatro áreas de conhecimento curricular e à educação profissional e técnica.
Nesta dimensão[42 ], tem -se um conjunto de três unidades curriculares, que, em diálogo, constituem -se nas itinerâncias, as quais as/os estudantes do Pará poderão criar condições para diversificar suas trajetórias e atender, em \textbf{certa} medida, suas expectativas e os interesses coletivos dos grupos juvenis que integram os diferentes estabelecimentos escolares de Ensino Médio no estado. \textbf{Figura} 6 : Esquematização das nucleações do currículo da nova arquitetura do ensino \textbf{médio} no Pará Fonte : Elaboração própria ( ProBNCC – etapa ensino \textbf{médio}, 2019 ).
É com base nesta organização curricular flexível por nucleação das áreas de conhecimento e das itinerâncias, conforme apresentado nas discussões das \textbf{figuras} 04 e 05, e sintetizadas na \textbf{figura} 06 acima, é que se estabeleceu a nova \textbf{matriz} curricular do Ensino Médio[43 ], disposta pedagogicamente em dois semestres por ano, com 500h de atividades de ensino, cada.
Assim, o semestre possuirá uma organização de 500h, distribuídas em 300h para a formação geral básica e 200h, destinadas à formação para o mundo do trabalho, totalizando 1.000h anuais e 3.000h mínimas, nos 03 anos do Ensino Médio.
A carga-horária total ( 3.000h mínimas ) deverá, ainda, ser organizada obrigatoriamente em 1.800h totais de formação geral básica e no mínimo 1.200h de formação para o mundo do trabalho, durante os 200 dias letivos, conforme estabelece a Lei nº 13.415/2017 em seu Art. 1º, Inciso I, § 1º[44 ], que alterou o Art. 24 da LDB/96.
Desta forma, a proposição genérica da nova arquitetura do ensino \textbf{médio} apresenta uma \textbf{matriz} curricular ajustada, de organização anual e funcionamento semestral, de acordo com as orientações da Lei nº 13.415/2017.
Nela, pode -se observar a distribuição da carga-horária, \textbf{tanto} por semestre, quanto por nucleação, a partir de suas unidades curriculares ( campos de saberes e práticas do ensino das áreas de conhecimento, presentes na formação geral básica ; e os projetos integrados e os campos de saberes e práticas eletivos, que se articulam às mesmas áreas de conhecimento e a educação profissional e técnica, \textbf{configurando}, portanto, a formação para o mundo do trabalho ).
Ressalta -se, nesta perspectiva, que há uma quarta unidade curricular, alocada na nucleação da formação para o mundo do trabalho e que tem por finalidade promover a integração entre as nucleações.
Trata -se do projeto de vida, como unidade curricular obrigatória no ensino \textbf{médio}.
As \textbf{figuras} 07 a 09, abaixo, apresentam a proposição da nova \textbf{matriz} curricular para o ensino \textbf{médio}, bem como o \textbf{desdobramento} do cálculo de distribuição de sua carga-horária e a exemplificação de um \textbf{quadro} horário de \textbf{aula}, para tornar mais clara a organização desta proposição. [ ^40 ] : A formação geral básica será tema da 4ª \textbf{seção} deste documento curricular, na qual se apresentará a concepção epistemológica e os \textbf{fundamentos} pedagógicos da área de conhecimento. [ ^41 ] : O detalhamento desta nucleação será apresentado, na \textbf{seção} 4, item 4.2, deste Documento Curricular.
Entretanto, vale ressaltar que há diferenças \textbf{conceituais} e ideológicas entre os termos “ mundo do trabalho ” e “ \textbf{mercado} de trabalho ” e estes \textbf{dizem} respeito aos conceitos de trabalho e emprego, respectivamente.
Assim, para o contexto desse documento, a formação para o mundo do trabalho visa ampliar a perspectiva da \textbf{categoria} trabalho, buscando compreender as relações \textbf{hegemônicas} de produção no sistema \textbf{capitalista} vigente.
Desta forma, considera -se importante uma formação que possibilite ao estudante do ensino \textbf{médio} compreender as relações entre os grupos que detêm os modos de produção e aqueles que disponibilizam suas forças de trabalho, para se entender, de forma mais ampla e \textbf{crítica}, as relações com o \textbf{mercado} e o emprego e a perspectiva ontológico-histórica da \textbf{categoria} trabalho. [ ^42 ] : É importante \textbf{frisar} que, como se trata de uma nucleação cuja finalidade é o \textbf{aprofundamento}, os objetos de conhecimento que serão trabalhados estão associados às quatro áreas de conhecimento e à educação profissional e técnica. [ ^43 ] : Trata -se de uma \textbf{matriz} curricular genérica para o ensino \textbf{médio} regular e que servirá de ponto de partida para a elaboração das \textbf{matrizes} curriculares para as modalidades de ensino e formas diversificadas de oferta do ensino \textbf{médio} no Pará.
Caberá a cada rede de ensino estabelecê -las, de acordo com o Conselho Estadual de Educação do Pará. [ ^44 ] : É importante ressaltar que este \textbf{parágrafo} traz a seguinte redação : “ § 1º A carga horária mínima anual de que trata o inciso I do caput deverá ser ampliada de forma progressiva, no ensino \textbf{médio}, para mil e quatrocentas ”.

% matched lemmas: abordagem, aluno, aprofundamento, ator, aula, autor, capacidade, capitalista, categoria, certo, ciência, competência, conceitual, configurar, contemporâneo, contextualização, crítico, debate, depender, desdobramento, detalhar, disciplina, dizer, econômico, figura, figurar, final, frisar, fundamento, grifo, hegemônico, ilustrar, imprimir, inerente, inter-relação, interdisciplinaridade, julgar, juventude, maneira, matriz, mercado, metodológico, mobilizar, médio, método, parágrafo, preocupar, problema, quadro, remeter, seção, sistematização, subsidiar, tanto, tecnologia, teoria, verificar
\end{textsample}
